% 6ページ版（../paper_asj6/main.tex、bb21a02）に、先輩の添削（2026-11-08、印刷への書き込み）をすべて反映した版
% 主な変更: はじめにから結果を除く、結果の章にあった方法を2.3〜2.5節へ、解釈を考察へ移す、VibraVox の説明を実験条件へ、
%           図1の文字を大きく（scripts/87 --large）、図2を2段幅にして差の段を追加（scripts/94 plot_wide）、疑似喉マイクを表に、
%           読点をカンマに、「であった」を「だった」に、括弧を減らす、略語のスペルアウト、式の記号を1文字の添字に、exp 表記
% ビルド: cd results/paper_asj7 && tectonic -X compile main.tex
% 数値の出典: results/GENERATIVE_RESULTS_2026-09-30.md §6・§7、results/review_response/stats_taps.csv・stats_vibravox.csv、results/acoustic/、results/jamo/
% 注意: 学会の正式テンプレートではなく、余白（上下18 mm・左右22 mm・段間6 mm）を合わせた自作の書式。投稿時は学会配布のスタイルファイルに載せ替える
\documentclass[xelatex,ja=standard,a4paper,twocolumn,10pt]{bxjsarticle}
\geometry{top=18mm,bottom=18mm,left=22mm,right=22mm,columnsep=6mm}
\setjamainfont[BoldFont={Hiragino Sans W6}]{Hiragino Mincho ProN}
\setjasansfont[BoldFont={Hiragino Sans W6}]{Hiragino Sans W3}
\usepackage{amsmath,amssymb}
\usepackage{graphicx}
\usepackage{booktabs}
\usepackage[font=small,skip=3pt]{caption}
\usepackage{cite}

\pagestyle{empty}
\setlength{\textfloatsep}{10pt plus 2pt minus 2pt}
\setlength{\floatsep}{6pt plus 2pt minus 2pt}

\makeatletter
\renewcommand{\section}{\@startsection{section}{1}{\z@}{1.0ex plus .3ex}{.5ex}{\normalfont\large\sffamily\bfseries}}
\renewcommand{\subsection}{\@startsection{subsection}{2}{\z@}{.7ex plus .2ex}{.3ex}{\normalfont\normalsize\sffamily\bfseries}}
\makeatother

\newcommand{\figdir}{../figures}
\setCJKfamilyfont{hangul}{AppleMyungjo}
\newcommand{\ko}[1]{{\CJKfamily{hangul}#1}}

\begin{document}

\twocolumn[{%
  \begin{center}
    {\Large\sffamily\bfseries 喉マイク音声の認識には，どんな前処理が効くのか\\[2pt]
      ——事前学習済み生成モデルの適応とWhisper用音声強調の役割——\footnotemark}\\[8pt]
    {春日　裕次（明治大学　総合数理学部　先端メディアサイエンス学科）}
  \end{center}
  \vspace{6pt}
}]
\footnotetext{What preprocessing helps throat-microphone speech recognition? Roles of pretrained generative restoration and Whisper-oriented speech enhancement, by KASUGA, Yuji (Meiji University).}

%────────────────────────────────────────
\section{はじめに}
喉マイクは頸部の振動を捉えるため環境騒音に強いが，高域や一部の子音の情報が乏しい。
韓国語の喉マイクと気導マイクの同時収録コーパスTAPS\cite{taps}の喉マイク音声を，本稿で損失の計算に用いるWhisper-small\cite{whisper}でそのまま認識すると，文字誤り率（character error rate: CER）は0.447になる。
同じ発話を気導マイクで収録した音声のCERは0.126であり，喉マイク音声ではその3倍以上になる。

音声を認識器に入力する前に変換する処理を，本稿では前処理と呼ぶ。
代表的な前処理は音声強調（speech enhancement: SE）である。
しかし，SEは信号の品質を改善しても，処理による歪みのために音声認識（automatic speech recognition: ASR）を悪化させうることが知られている\cite{ochiai2024}。
喉マイクでも，汎用の帯域拡張や声質変換がCERを悪化させたと報告されている\cite{yamanaka2026}。この報告では，TAPSで追加学習したモデルでも悪化し，CERは正解テキストではなく気導音の認識結果を基準に計算されている。
一方，Hauretら\cite{hauret2025rt}は，Kuら\cite{ku2024}の事前学習済みの音声復元モデルを仏語のコーパスVibraVox\cite{vibravox}の喉マイク音声で追加学習し，1種の音素認識器で誤りが減ったと報告した。本稿もこの音声復元モデルを用いる。
また，凍結したASRの損失で前処理のモデルを学習する方法\cite{dissen2024}や，生成型と予測型のSEの出力を組み合わせる方法\cite{shi2024}も提案されている。

本稿の目的は，喉マイク音声に対して前処理が認識を助ける条件と，ASRの損失で学習したSEの出力を生成モデルの出力に組み合わせたときに加わる改善，およびその改善に関わる信号の成分を調べることである。
組み合わせ方の新しい提案は目的としない。
評価には，入力の特徴量や学習方法が異なる複数の認識器を用い，TAPSに加えてVibraVoxでも同じ手順を繰り返す。

%────────────────────────────────────────
\section{方法}
\begin{figure}[t]
  \centering
  \includegraphics[width=\columnwidth]{\figdir/method_overview_large.png}
  \caption{手法の全体像。喉マイク音声を生成モデルとWhisper用SEの両方に通し，出力の振幅を融合して9種の認識器で評価する}
  \label{fig:overview}
\end{figure}

\subsection{生成モデルによる復元}
\label{sec:gen}
手法の全体像を図\ref{fig:overview}に示す。
NVIDIAが公開している，flow matching\cite{lipman2023}で学習された音声復元モデル\cite{ku2024}を用いる。
このモデルは，複素短時間フーリエ変換（short-time Fourier transform: STFT）の係数を扱う24層のTransformerで，パラメータ数は4.3億である。
英語約6万時間の音声で，マスクした区間を復元するタスクにより事前学習されている。

喉マイクと気導マイクで同時に収録した音声の組を，以下では対データと呼ぶ。
本稿では，喉マイク音声を条件，同時に収録した気導マイク音声を目標として，対データでこのモデルを追加学習する。
この追加学習を適応と呼ぶ。適応ではASRの損失を用いない。
推論では標準正規分布の雑音から出発し，flow matchingの常微分方程式をEuler法の20ステップで解いて生成する。
生成結果は出発点の雑音によって変わるので，$N=4$個のサンプルのSTFT振幅を平均し，位相には1つ目のサンプルのものを用いる。
この出力を生成×4と書く。

\subsection{Whisper用SEと振幅の融合}
\label{sec:fusion}
TAPS論文\cite{taps}が公開しているSE-Conformer\cite{seconformer}を，以下ではTAPS SEと呼ぶ。
TAPS SEは，TAPSのtrainで喉マイク音声を気導マイク音声に変換するように，再構成損失で学習されている。
これを初期値として，次の損失$L$で追加学習したSEをCE-SEと呼ぶ。
\begin{equation}
  L = L_\mathrm{r}(\hat{s}, a) + \lambda\, L_\mathrm{c}(\hat{s}, y)
  \label{eq:loss}
\end{equation}
ここで，$\hat{s}$はSEの出力，$a$は気導マイク音声，$y$は正解テキスト，$\lambda$は重みである。

$L_\mathrm{r}$は再構成損失で，TAPS SEの学習\cite{taps}と同じく，波形のL1損失と多重解像度STFT損失\cite{pwg}の和とした。
多重解像度STFT損失は，FFT長が512，1024，2048点の3つの解像度で，振幅のspectral convergenceと対数振幅のL1誤差を0.5ずつ足し，解像度について平均したものである。
係数0.5もTAPS SEの学習と同じ値である。

$L_\mathrm{c}$は交差エントロピー（cross entropy: CE）の損失である。
$\hat{s}$からWhisperと同じ対数メルスペクトログラムを微分可能な形で求めてWhisper-small\cite{whisper}に入力し，正解テキストのトークン列をデコーダに与えて，各トークンの予測の交差エントロピーを求める。言語は韓国語に指定する。
Whisper-smallの重みは更新せず，$L_\mathrm{c}$の勾配をWhisperを通してSEに伝える。
学習のたびにWhisperを通して勾配を計算するので，計算量を抑えるため，Whisperの中で比較的小さいsmallを用いた。

次に，生成×4とCE-SEの出力を組み合わせる。
両者のSTFT（512点，シフト128点，Hann窓）の複素係数をそれぞれ$G$，$C$とし，次式で振幅を重み付きで平均し，位相には$G$の位相を用いる。
\begin{equation}
  Y = \bigl\{(1-w)\,|G| + w\,|C|\bigr\}\, \exp(j\angle G)
  \label{eq:fusion}
\end{equation}
$Y$は組み合わせた後の複素係数，$j$は虚数単位である。
本稿では，この操作を振幅の融合，または単に融合と呼ぶ。
$\lambda=10$は$\{0, 0.1, 0.5, 1, 2, 5, 10\}$から，CE-SE側の重み$w=0.5$は生成1回の出力との融合で$\{0.25, 0.5, 0.75\}$から，いずれもdevのCERで選んだ。

\subsection{疑似喉マイク音声}
\label{sec:pseudo}
喉マイクの認識の難しさが，平均的な周波数特性と帯域の狭さで説明できるかを調べるため，気導マイク音声から疑似的な喉マイク音声を作る。以下ではこれを疑似喉マイクと呼ぶ。
TAPS trainの最大1,000対から，気導マイク音声のエネルギーが発話内の最大から40~dB以内のフレームを選ぶ。
それらのフレームで，喉マイクと気導マイクのパワーを周波数ビン$f$ごとに合計し，それぞれ$S_\mathrm{t}(f)$，$S_\mathrm{a}(f)$とする。
TAPS testの気導マイク音声のSTFT（512点，シフト128点）に振幅比$H(f)=\sqrt{S_\mathrm{t}(f)/S_\mathrm{a}(f)}$をかけ，4~kHz以上の係数を0にする。
最後に，RMSをtrainの喉マイク音声のRMSの中央値にそろえる。

\subsection{字母誤り率}
\label{sec:jamo}
韓国語のどの音の認識誤りが減ったかを見るため，ハングルの字母，つまり子音字と母音字の単位で誤りを数える。
ハングルの音節1文字は，初声の子音字，中声の母音字，終声の子音字からなる。終声は無い場合もある。
Unicodeのハングル音節は，符号位置から3つの字母に算術的に分解できる。
そこで，句読点と空白を除いた認識結果と正解を，それぞれ字母の列に分解する。
両者を単位コストの編集距離で対応づけ，正解の字母のうち置換か脱落になったものの割合を，字母の種類ごとに全発話で合算して求める。
これを字母誤り率と呼ぶ。挿入は数えない。
この率は正書法上の字母の誤りで，実際に発音された音の誤りとは一致しない。

字母の種類は音声学の区分に従い，母音，鼻音~\ko{ㄴㅁ}，平音~\ko{ㄱㄷㅂㅈ}，激音~\ko{ㅋㅌㅍㅊ}，摩擦音~\ko{ㅅㅆㅎ}に分ける。
激音は強い息を伴う子音で，平音は同じ調音位置で息の弱い子音である。平音と激音には，破裂音と破擦音の初声を数える。
正解は数をハングルで書くが，認識結果では数字になることがある。このとき\ko{칠}（7）や\ko{팔}（8），\ko{천}（千）に含まれる激音の誤りが水増しされる。
これを避けるため，各認識器について，処理なし，TAPS SE，CE-SE，生成×4，融合，気導の6条件のどの認識結果にも数字とラテン文字が出なかった発話だけを使う。
この選び方は認識結果に依存するので，選んだ発話はtest全体を代表しない。

処理なしと気導の差のうち，生成×4で残った割合を$r=(e_\mathrm{g}-e_\mathrm{a})/(e_\mathrm{n}-e_\mathrm{a})$とする。
$e_\mathrm{n}$，$e_\mathrm{g}$，$e_\mathrm{a}$は，それぞれ処理なし，生成×4，気導の字母誤り率である。$r$は話者ごとに求めて中央値をとる。
激音と平音の$r$の比較は，結果を見てから選んだ。
そのためHolm補正は3種の認識器の間の多重比較だけを補正し，比較の選び方による偏りは補正していない。この分析は探索的である。

\subsection{融合の対照と成分の切り分け}
\label{sec:decomp}
融合による改善が単に2つの出力を平均した効果かを調べるため，CE-SEの代わりに，CE損失を用いないSEを同じ方法で生成×4と融合する。
用いるSEは，TAPS SE，$L_\mathrm{r}$だけで追加学習したSE（$\lambda=0$），SSL-MSEである。
SSL-MSEは，Satoら\cite{sato2025}の自己教師あり学習（self-supervised learning: SSL）表現の損失を参考にした対照である。
WavLM-Large\cite{wavlm}の全層で表現の平均二乗誤差（mean squared error: MSE）を求めて平均し，重み0.3で$L_\mathrm{r}$に加えてTAPS SEを追加学習した。
重み0.3は，SSLの項が$L_\mathrm{r}$の約13〜14倍になるように決めた。これはCE-SEの$\lambda=10$での$L_\mathrm{c}$の項と同じ比である。学習は最大15エポックで，devの損失で重みを選んだ。
Satoらは後半の層を重み付きで平均した表現の誤差を用いており，本稿の実装とは異なる。

次に，融合が生成×4から何を変えたかを，対数振幅の差$D=\log|Y|-\log|G|$で調べる。
$G$と$Y$は\ref{sec:fusion}節と同じSTFTで求め，振幅に$10^{-7}$を足してから自然対数をとる。
フレームの種類は，気導マイク音声を使わずに生成×4の出力から決める。
パワーが発話内の最大から40~dB以上低いフレームを無音，無音でなく2~kHz未満のパワーの割合が0.5を超えるフレームを有声らしいフレーム，残りを無声とする。

第1の操作では，$D$を無音以外のフレームで時間平均した周波数ごとの値を，生成×4の対数振幅に加える。これを平均ゲインと呼ぶ。
平均ゲインは，devで求めたもの，フレームの種類ごとに求めたもの，同じ発話の融合から求めたものの3通りとする。最後のものは発話ごとのイコライザに当たり，発話別EQと呼ぶ。

第2の操作では，$D$を包絡と細部に分ける。
各フレームの$D$（257点）を逆フーリエ変換して512点の実ケプストラムとし，ケフレンシー0〜29点とそれに対称な点だけを残してフーリエ変換し，実部をとる。
この操作を$P$，恒等写像を$I$とし，$PD$を包絡，$(I-P)D$を細部と呼ぶ。
29点は16~kHzで約1.9~msで，ピッチ周期の2.5~ms以上より短い。この次数は結果を見る前に決めた。
包絡か細部の片方だけを，フレームごとに生成×4の対数振幅に加える。両方を加えると融合の対数振幅に戻る。位相は生成×4のものに固定する。

第3の操作では，有声らしいフレームでだけ融合の出力を使い，他のフレームでは生成×4を使う。
比較として，同じマスクを発話ごとに乱数で選んだ量だけ時間方向に巡回シフトしたものも作る。マスクは3フレームの移動平均でなめらかにする。
有声らしいフレームは全体の約83\%で，無音以外のフレームでは，気導マイク音声に対するpYIN\cite{pyin}の有声判定との一致率が0.88だった。

最後に，融合による音響特徴量の変化を測る。
高域のずれは，4~kHz以上の対数メルの，気導マイク音声との平均絶対差とする。
スペクトル傾斜は，気導マイク音声のpYINで有声と判定した区間で，4~kHz未満の対数メルに当てはめた回帰直線の傾きとする。
フォルマントの帯域幅は線形予測（LPC）で推定する。同じ有声区間で，信号を8~kHzに下げてプリエンファシス（係数0.97）をかけ，25~msのハミング窓（10~msシフト）で10次のLPCを求め，90〜3800~Hzにあり帯域幅が600~Hz未満の極を周波数の低い順に3つとる。
これらの量の発話ごとの変化とCERの変化の関係を，スピアマン順位相関係数$\rho$で調べる。

%────────────────────────────────────────
\section{実験条件}
\label{sec:setup}
TAPS\cite{taps}の話者独立な分割を用いた。trainは40話者，devとtestは10話者ずつで，各話者100発話である。
testの話者は学習と重ならないが，収録機器と言語は同じである。
音声はすべて16~kHzで扱った。TAPSの喉マイク音声は8~kHzで収録された。

別の喉マイクとして，仏語の体内伝導マイクコーパスVibraVox\cite{vibravox}のspeech\_cleanを用いた。speech\_cleanは静かな環境での読み上げ音声で，喉マイク音声と，同時に収録したヘッドセットの気導マイク音声を使う。
trainは20,981対，testは21話者の3,064発話で，testの話者はtrainと重ならない。
testの文のうちtrainにも現れるのは1文だけだった。TAPSでは1,000文中6文だった。

生成モデルはTAPS trainの4,000対で2万回更新した。1バッチは6.14秒の切り出し8個，学習率は$10^{-4}$で，最終ステップの重みを用いた。
CE-SEは15秒を超える発話を除き，バッチ4発話，学習率$3\times10^{-4}$で最大50エポック学習した。devでの式\eqref{eq:loss}の損失が5エポック下がらなければ止め，損失が最小のエポックの重みを用いた。
VibraVoxでも同じ手順で適応した。生成モデルは事前学習済みの重みから適応した。SEはTAPS SEを$L_\mathrm{r}$だけでVibraVoxに合わせてから，仏語の正解テキストでCE-SEを学習した。$\lambda$，$w$，$N$はTAPSの値に固定した。
VibraVoxでは，trainの1,000対で1万回更新した適応も1回だけ試した。

CE-SEと生成モデルは，最初に学習したもの（以下，元の学習）に加えて，乱数のシードと数値精度を含む計算環境を変えて2回ずつ学習し直した。これらを学習し直し1，2と呼ぶ。
同じ回どうしを融合し，3回とも同じtestの10話者で評価した。

Whisper\cite{whisper}と同じ形式の対数メルを入力とする次の6種を，本稿ではWhisper系と呼ぶ。Whisper-small，base，medium，large-v3-turbo，Qwen3-ASR-1.7B\cite{qwen3asr}，FT Whisperである。
FT Whisperは，TAPS trainの喉マイク音声と正解テキストで，Whisper-smallの全パラメータを追加学習したモデルである。学習率は$10^{-5}$，バッチは16で，devのCERが最小のエポックを用いた。
これらと入力の特徴量や学習方法が異なる3種を，系統の異なる認識器と呼ぶ。
波形を入力とし，connectionist temporal classification（CTC）で学習されたMMS-1B\cite{mms}とXLSR，および韓国語の自発音声コーパスKsponSpeechで学習されたZipformer\cite{zipformer}である。
MMS-1Bには韓国語のアダプタを用いた。XLSRは，XLSR-53\cite{xlsr}を韓国語でfine-tuningした公開モデルである。Zipformerはtransducer型の公開モデルである。
Whisper系は言語を指定した。Whisperはビーム幅5，CTCの2種とZipformerは貪欲法で復号した。
VibraVoxでは，韓国語だけのFT WhisperとZipformerを除き，XLSRとMMSを仏語版に替えた7種を用いた。

CERは文字単位で数えた。韓国語ではハングルの音節1文字を1文字とする。
句読点を除き，仏語は小文字にし，空白も1文字として数えた。発話ごとのCERを1.0で打ち切ってから全発話で平均した。
全発話の編集距離の合計を正解の文字数の合計で割る通常のCERでも，主な傾向は同じだった。
Zipformerは空白を出力しないため，この数え方ではCERが高く出るが，空白を除いたCERでも主な条件の大小関係は同じだった。

話者ごとの平均CERで，両側のWilcoxon符号順位検定を行った。
条件の対ごとに，その対を評価したすべての認識器のp値をまとめてHolm法で補正し，補正後の$p<0.05$を有意とした。この補正は全探索に対するものではない。
TAPSの10話者では，9種で補正した後のpの最小値は0.018である。
融合を採用するかは，Whisper-small以外の認識器のtestの結果も見て決めた。このため，TAPS上の分析は探索的である。

%────────────────────────────────────────
\section{実験結果}
\subsection{前処理が効く条件}
\label{sec:res_main}
表\ref{tab:main}に各手法のCERを示す。
生成×4は9種すべてでTAPS SEよりCERを下げた。$N=1$の1回の生成でもTAPS SEより15〜24\%低く，10話者全員で下がった。
FT Whisperでも，処理なしの0.138から，生成×4で0.100，融合で0.092へ下がった。融合は処理なしより33\%低く，10話者全員で下がった。

同じ構造のモデルを，同じ2万回の更新だけ乱数の初期値から学習すると，Whisper-smallのCERは0.551だった。
VibraVoxで適応した生成モデルをTAPSにそのまま用いると，1回の生成で0.628だった。
どちらも処理なしの0.447より高い。後者と同時に認識し直した処理なしの値は0.448だった。
融合後も，Whisper-smallで気導マイク音声の0.126との差は残った。

\begin{table}[t]
  \centering
  \caption{TAPS testのCER。太字は喉マイク入力の条件内での最小。なし: 処理なし，TAPS: TAPS SE，W-FT: FT Whisper，気導: 同じ発話の気導マイク音声をそのまま認識した参考値で3種のみ測定，*: 損失・選択に使用}
  \label{tab:main}
  \footnotesize
  \setlength{\tabcolsep}{3.0pt}
  \begin{tabular}{lcccccc}
    \toprule
    認識器 & なし & TAPS & CE-SE & 生成×4 & 融合 & 気導 \\
    \midrule
    W-small* & 0.447 & 0.230 & 0.196 & 0.169 & \textbf{0.163} & 0.126 \\
    W-base & 0.655 & 0.273 & 0.267 & 0.205 & \textbf{0.201} & -- \\
    W-medium & 0.337 & 0.196 & 0.169 & 0.150 & \textbf{0.144} & -- \\
    W-turbo & 0.257 & 0.170 & 0.141 & 0.131 & \textbf{0.125} & 0.103 \\
    W-FT & 0.138 & 0.146 & 0.121 & 0.100 & \textbf{0.092} & -- \\
    Qwen3 & 0.243 & 0.136 & 0.114 & 0.097 & \textbf{0.094} & 0.072 \\
    \midrule
    XLSR & 0.584 & 0.222 & 0.279 & \textbf{0.161} & 0.163 & -- \\
    MMS-1B & 0.497 & 0.350 & 0.354 & \textbf{0.284} & 0.285 & -- \\
    Zipformer & 0.997 & 0.542 & 0.791 & \textbf{0.433} & 0.483 & -- \\
    \bottomrule
  \end{tabular}
\end{table}

\begin{table}[t]
  \centering
  \caption{疑似喉マイクのCER（TAPS test）。喉マイク: 本物の喉マイク音声で，表\ref{tab:main}とは別の回の認識結果}
  \label{tab:pseudo}
  \footnotesize
  \setlength{\tabcolsep}{4.0pt}
  \begin{tabular}{lccc}
    \toprule
    認識器 & 疑似喉マイク & 気導 & 喉マイク \\
    \midrule
    W-small & 0.130 & 0.126 & 0.448 \\
    W-turbo & 0.103 & 0.103 & 0.258 \\
    Qwen3 & 0.077 & 0.072 & 0.243 \\
    \bottomrule
  \end{tabular}
\end{table}

\subsection{疑似喉マイクと字母ごとの誤り}
表\ref{tab:pseudo}に疑似喉マイクのCERを示す。
Whisper-smallでは，疑似喉マイクと気導マイク音声のCERの差は0.003で，本物の喉マイクとの差よりはるかに小さかった。

\begin{figure*}[t]
  \centering
  \includegraphics[width=\textwidth]{\figdir/spectrogram_gen_fusion_wide.png}
  \caption{TAPS testの1発話（話者p00）の冒頭4秒の振幅スペクトログラム。上から喉マイク（処理なし），生成×4，融合，融合と生成×4の差（dB），気導マイク。破線は4~kHz。差の段では，両方の振幅が各図の最大から50~dB以上低いビンを0として表示した}
  \label{fig:spec}
\end{figure*}

図\ref{fig:spec}に1発話の振幅スペクトログラムを示す。
生成×4は喉マイクに無い4~kHz以上の成分を補い，その時間的な配置は多くの区間で気導マイクと対応して見えた。喉マイクの無音区間にある小さな雑音は，生成×4では消えた。
融合と生成×4は上の段どうしを比べると違いが目立たないが，差の段では，一部の区間の高域で融合が振幅を下げ，別の区間では上げており，変化の正負は時間と周波数によって入れ替わっていた。

字母誤り率を求めた発話は，Whisper-smallで309，turboで340，Qwen3で450発話で，いずれも10話者を含む。
表\ref{tab:jamo}にWhisper-smallの字母誤り率を示す。
処理なしでは激音の字母誤り率が0.423と最も高く，平音の0.277を上回った。
Whisper-smallでは，生成×4はすべての種類で字母誤り率を処理なしの約4分の1以下に下げた。
$r$は激音で18\%，平音で10\%だった。
激音の方が$r$が大きく，Whisper-smallとturboでは有意だった。どちらも10話者中8話者で激音の$r$の方が大きかった。Qwen3では有意でなかった。
対応づけを逆方向から求める，3種に共通の283発話に限る，全1,000発話で集計する，のいずれでも，Whisper-smallの差は有意のまま残ったが，turboの差は有意でなくなった。
摩擦音も$r$が大きい傾向があったが，破裂音との差は有意でなかった。

\begin{table}[t]
  \centering
  \caption{字母誤り率（Whisper-small，309発話）。残り: $r$の話者ごとの中央値で，10話者とも分母は正}
  \label{tab:jamo}
  \footnotesize
  \setlength{\tabcolsep}{3.0pt}
  \begin{tabular}{lrcccc}
    \toprule
    種類 & 字母数 & なし & 生成×4 & 気導 & 残り \\
    \midrule
    母音 & 14104 & 0.282 & 0.050 & 0.021 & 13\% \\
    鼻音 & 1151 & 0.321 & 0.045 & 0.017 & 9\% \\
    平音 & 5112 & 0.277 & 0.034 & 0.010 & 10\% \\
    激音 & 1067 & 0.423 & 0.094 & 0.023 & 18\% \\
    摩擦音 & 2712 & 0.225 & 0.054 & 0.021 & 19\% \\
    \bottomrule
  \end{tabular}
\end{table}

\subsection{ASR損失で学習したSEとの融合}
融合は，Whisper系6種すべてで生成×4よりCERを下げた（$-1.8$〜$-8.0$\%）。以下，この追加の改善を上乗せと呼ぶ。
各学習回のCERと補正後のp値を表\ref{tab:stat}に示す。
small，medium，turbo，FTでは，元の学習を含む3回とも有意だった（補正後$p\le0.020$）。baseは学習し直し1だけで有意で，Qwen3はどの回も有意でなかった。
CE-SE単体は，元の学習ではbase以外のWhisper系5種をTAPS SEより有意に改善したが，XLSRとZipformerはTAPS SEより悪化させた。
悪化は元の学習で$+26$\%と$+46$\%，学習し直した2回では約$+11$〜$12$\%だった。処理なしと比べると，どちらも改善した。

\begin{table*}[t]
  \centering
  \caption{主要な比較の各学習回のCERと，話者単位のWilcoxon検定のHolm補正後のp値（TAPS test，3回とも同じ10話者）。太字は$p<0.05$。融合のp: 生成×4との比較。包絡のp: 生成×4との比較と，融合との比較。包絡は元の学習のみ}
  \label{tab:stat}
  \footnotesize
  \setlength{\tabcolsep}{3.0pt}
  \begin{tabular}{l ccc ccc ccc ccc}
    \toprule
    & \multicolumn{3}{c}{元の学習} & \multicolumn{3}{c}{学習し直し1} & \multicolumn{3}{c}{学習し直し2} & \multicolumn{3}{c}{包絡（元の学習）} \\
    \cmidrule(lr){2-4}\cmidrule(lr){5-7}\cmidrule(lr){8-10}\cmidrule(lr){11-13}
    認識器 & 生成×4 & 融合 & p & 生成×4 & 融合 & p & 生成×4 & 融合 & p & CER & p（対生成×4） & p（対融合） \\
    \midrule
    W-small* & 0.169 & 0.163 & \textbf{0.020} & 0.169 & 0.162 & \textbf{0.018} & 0.169 & 0.162 & \textbf{0.018} & 0.166 & 0.078 & 0.059 \\
    W-base & 0.205 & 0.201 & 0.422 & 0.205 & 0.201 & \textbf{0.029} & 0.204 & 0.200 & 0.137 & 0.204 & 0.984 & 0.193 \\
    W-medium & 0.150 & 0.144 & \textbf{0.018} & 0.151 & 0.143 & \textbf{0.018} & 0.150 & 0.142 & \textbf{0.018} & 0.146 & \textbf{0.049} & 0.193 \\
    W-turbo & 0.131 & 0.125 & \textbf{0.018} & 0.132 & 0.125 & \textbf{0.018} & 0.132 & 0.124 & \textbf{0.018} & 0.126 & \textbf{0.018} & 0.232 \\
    W-FT & 0.100 & 0.092 & \textbf{0.018} & 0.100 & 0.091 & \textbf{0.018} & 0.099 & 0.091 & \textbf{0.018} & 0.096 & \textbf{0.027} & 0.068 \\
    Qwen3 & 0.097 & 0.094 & 0.422 & 0.100 & 0.097 & 0.480 & 0.099 & 0.099 & 1.000 & 0.097 & 1.000 & \textbf{0.031} \\
    \midrule
    XLSR & 0.161 & 0.163 & 0.422 & 0.163 & 0.164 & 0.480 & 0.162 & 0.163 & 1.000 & 0.168 & \textbf{0.027} & \textbf{0.018} \\
    MMS-1B & 0.284 & 0.285 & 0.625 & 0.285 & 0.283 & 0.336 & 0.284 & 0.284 & 1.000 & 0.287 & 0.252 & 0.109 \\
    Zipformer & 0.433 & 0.483 & \textbf{0.018} & 0.432 & 0.437 & 0.770 & 0.431 & 0.438 & 0.336 & 0.510 & \textbf{0.018} & \textbf{0.031} \\
    \bottomrule
  \end{tabular}
\end{table*}

表\ref{tab:ctrl}に，CE損失を用いないSEとの融合の結果を示す。
TAPS SEや$\lambda=0$のSEとの融合は，Whisper系で生成×4から$-1.7$〜$+3.1$\%で，ほぼ改善しなかった。
CE-SEとの融合のCERは，これら2つの融合より6種すべてで有意に低かった。
SSL-MSEとの融合でも，turboとFTでは生成×4より有意にCERが下がった。

\begin{table}[t]
  \centering
  \caption{生成×4に融合するSEを変えた対照（元の学習，CER）。+TAPS: TAPS SE，+$\lambda$=0: $L_\mathrm{r}$だけで学習したSE，+SSL: SSL-MSE，+CE: CE-SEとの融合。baseは本文のみ}
  \label{tab:ctrl}
  \footnotesize
  \setlength{\tabcolsep}{3.0pt}
  \begin{tabular}{lccccc}
    \toprule
    認識器 & 生成×4 & +TAPS & +$\lambda$=0 & +SSL & +CE \\
    \midrule
    W-small* & 0.169 & 0.172 & 0.172 & 0.168 & \textbf{0.163} \\
    W-medium & 0.150 & 0.150 & 0.151 & 0.147 & \textbf{0.144} \\
    W-turbo & 0.131 & 0.130 & 0.129 & 0.128 & \textbf{0.125} \\
    W-FT & 0.100 & 0.099 & 0.098 & 0.095 & \textbf{0.092} \\
    Qwen3 & 0.097 & 0.099 & 0.100 & 0.096 & \textbf{0.094} \\
    \bottomrule
  \end{tabular}
\end{table}

\subsection{上乗せに関わる成分}
\label{sec:res_decomp}
表\ref{tab:why}に，元の学習での切り分けの結果を示す。
平均ゲインを加えたものは，3通りのいずれも，Whisper系で生成×4から$\pm2$\%以内で，どれも有意でなかった。
細部だけを加えても，どの認識器でも有意な改善はなかった。
包絡だけを加えると，medium，turbo，FTで生成×4より約3〜4\%下がって有意だった（表\ref{tab:stat}）。その量は，融合の改善量の約半分〜9割だった。この割合は点推定で，区間は求めていない。
この3種では，包絡だけと融合の差は有意でなかった。Qwen3では包絡だけでは改善せず，融合は包絡だけより有意に低かった。

有声らしいフレームでだけ融合を使うと，small，medium，turbo，FT，Qwen3で生成×4より有意に下がった。
マスクを巡回シフトしても，medium，turbo，FTで有意に下がった。有声らしいフレームと巡回シフトの差は，どの認識器でも有意でなかった。

融合は，高域のずれとスペクトル傾斜を気導マイク音声に近づけ，フォルマントの帯域幅を広げた。
Whisper-small，medium，FT，Qwen3では，これらの量の発話ごとの変化とCERの変化の$|\rho|$は約0.08以下だった。Zipformerでは約0.2の量もあった。

系統の異なる認識器では，Zipformerは発話別EQだけで生成×4より3.6\%改善し，10話者全員で下がった。
一方，包絡だけを加えると$+17.7$\%で，融合の$+11.7$\%より悪化した。XLSRも包絡だけで$+3.9$\%悪化した。

\begin{table}[t]
  \centering
  \caption{上乗せの切り分け（元の学習，CER）。EQ: 発話別EQ，包絡／細部: 融合と生成×4の対数振幅の差の包絡／細部だけを加えたもの}
  \label{tab:why}
  \footnotesize
  \setlength{\tabcolsep}{3.0pt}
  \begin{tabular}{lccccc}
    \toprule
    認識器 & 生成×4 & EQ & 包絡 & 細部 & 融合 \\
    \midrule
    W-small* & 0.169 & 0.170 & 0.166 & 0.169 & 0.163 \\
    W-medium & 0.150 & 0.151 & 0.146 & 0.151 & 0.144 \\
    W-turbo & 0.131 & 0.131 & 0.126 & 0.132 & 0.125 \\
    W-FT & 0.100 & 0.100 & 0.096 & 0.099 & 0.092 \\
    Qwen3 & 0.097 & 0.096 & 0.097 & 0.099 & 0.094 \\
    \midrule
    Zipformer & 0.433 & 0.417 & 0.510 & 0.444 & 0.483 \\
    \bottomrule
  \end{tabular}
\end{table}

\subsection{VibraVoxでの評価}
表\ref{tab:vbx}にVibraVox testのCERを示す。
気導マイク音声のCERは，devでWhisper-smallが0.032，turboが0.016，Qwen3が0.014だった。
TAPSで学習した生成モデルとSEをそのまま用いると，いずれも処理なしより悪化した。音量と帯域をTAPSにそろえても，処理なしより悪いままだった。

VibraVoxで適応し直すと，7種の認識器で，生成×4は処理なしよりCERを70〜82\%下げ，7種とも補正後$p<0.001$だった。
融合は生成×4より，点推定でさらに3〜15\%CERを下げ，base以外の6種で有意だった（$p\le0.007$）。baseは$p=0.055$だった。
CE-SEと融合は，TAPSと異なり，仏語のXLSRとMMSも改善した。
Whisper-smallと1回の生成では，trainの1,000対，約1.3時間での適応でも，全データ約26時間での改善量の約86\%が得られた。

\begin{table}[t]
  \centering
  \caption{VibraVox testのCER（仏語，7種）。太字は各列の最小。上段: TAPSで学習したまま，下段: VibraVoxで適応。再構成のみ: $L_\mathrm{r}$だけでVibraVoxに合わせたTAPS SE。W-s・W-b・W-m・W-t: Whisper-small・base・medium・turbo。XLSR・MMSは仏語版。--は未評価}
  \label{tab:vbx}
  \footnotesize
  \setlength{\tabcolsep}{1.6pt}
  \begin{tabular}{lccccccc}
    \toprule
    条件 & W-s & W-b & W-m & W-t & Qwen3 & XLSR & MMS \\
    \midrule
    処理なし & 0.422 & 0.536 & 0.366 & 0.269 & 0.242 & 0.275 & 0.258 \\
    TAPS SE & 0.711 & 0.755 & 0.718 & 0.654 & 0.594 & 0.501 & 0.577 \\
    生成×4 & 0.559 & 0.594 & 0.542 & 0.497 & 0.454 & -- & -- \\
    \midrule
    再構成のみ & 0.256 & 0.313 & 0.212 & 0.163 & 0.154 & 0.163 & 0.187 \\
    CE-SE & 0.153 & 0.226 & 0.114 & 0.081 & 0.081 & 0.128 & 0.131 \\
    生成×4 & 0.087 & 0.123 & 0.067 & 0.053 & 0.046 & 0.063 & 0.076 \\
    融合 & \textbf{0.079} & \textbf{0.119} & \textbf{0.059} & \textbf{0.045} & \textbf{0.041} & \textbf{0.058} & \textbf{0.067} \\
    \bottomrule
  \end{tabular}
\end{table}

\subsection{計算量}
融合は4.3億パラメータのモデルを20ステップ，4サンプル実行するため重い。
表\ref{tab:steps}に，Euler法のステップ数とサンプル数を減らした結果を示す。
5ステップ，1サンプルでは，生成モデルの推論時間が1発話あたり約4.1秒から約0.41秒へ，約1/10になった。平均発話長は約9秒である。
この時間はモデルの読み込みを含む一括処理の時間で，逐次処理の遅延は測っていない。
20×4に対する5×1の融合のCERの変化は，Whisper系で$-0.03$〜$+3.6$\%で，補正後に有意に悪化したのはbaseの$+2.3$\%（$p=0.041$）だった。
系統の異なる認識器では，XLSRの$+4.3$\%（$p=0.018$）とZipformerの$+11.0$\%（$p=0.018$）が有意に悪化した。
5×4でも，XLSR（$p=0.047$）とZipformer（$p=0.018$）は20×4より有意に悪化した。

\begin{table}[t]
  \centering
  \caption{融合のステップ数・サンプル数とCER（元の学習，test 1,000発話）。条件: Euler法のステップ数×サンプル数。時間: 同じ100発話での生成モデルの推論時間で，RTX PRO 6000でモデルの読み込みを含む一括の時間。p: 20×4と5×1の話者単位のWilcoxon検定のHolm補正後のp値で，9種で補正。W-s・W-b・W-m: Whisper-small・base・medium，Zip.: Zipformer}
  \label{tab:steps}
  \footnotesize
  \setlength{\tabcolsep}{1.6pt}
  \begin{tabular}{lccccccc}
    \toprule
    条件 & 時間 & W-s & W-b & W-m & W-FT & XLSR & Zip. \\
    \midrule
    20×4（本稿） & 405 s & 0.163 & 0.201 & 0.144 & 0.092 & 0.163 & 0.483 \\
    5×4 & 115 s & 0.162 & 0.202 & 0.142 & 0.091 & 0.167 & 0.522 \\
    5×1 & 41 s & 0.164 & 0.206 & 0.144 & 0.095 & 0.170 & 0.536 \\
    p & & 1.000 & \textbf{0.041} & 1.000 & 0.186 & \textbf{0.018} & \textbf{0.018} \\
    \bottomrule
  \end{tabular}
\end{table}

%────────────────────────────────────────
\section{考察と限界}
前処理が喉マイク音声の認識を助けるかは，条件で反転した。
今回比較した条件では，対象の喉マイクで適応した事前学習済みの生成モデルは，認識器の系統によらず効いた。
同じ2万回の更新で乱数初期値から学習したときや，別の喉マイクで適応したままのときは，害になった。
ただし，事前学習済みのモデルを適応せずに用いる条件は評価しておらず，事前学習と適応のそれぞれが必要かは示していない。
また，TAPSの約10時間で4.3億パラメータのモデルを乱数から学習しきれていない可能性も残る。1210万パラメータのTAPS SEは，同じデータで乱数から学習されている。
したがって乱数初期値との比較が示すのは，同じ2万回の更新では，事前学習済みの重みから始めたときにだけ，この規模のモデルが約10時間の対データで有効になったことである。
先行研究\cite{yamanaka2026}でCERが悪化した手法との違いが事前学習の有無にあるかは，本稿では確かめていない。

喉マイクの認識の難しさは，今回の平均的な振幅特性と帯域制限による疑似喉マイクでは再現されなかった。時間とともに変わる，音によって異なる劣化が関わっている可能性がある。
Whisper-smallでの探索的な分析では，生成後も激音の字母の誤りが相対的に多く残った。
候補の一つは，激音に伴う強い気息，つまり声帯の振動を伴わない息の雑音の手がかりが，喉マイクで弱まることである。
ただし本分析は認識文字列に基づく。発声開始時間（voice onset time: VOT）や後続母音の基本周波数，表記と発音の違いを調べていないため，原因は特定できない。

CE損失による学習は上乗せを大きくした。しかし，CE損失を用いないSSL-MSEとの融合でも一部の認識器で上乗せが得られたので，CE損失は上乗せの必要条件ではない。
元の学習では，上乗せは検討した平均ゲインでは再現せず，融合と生成×4の差の包絡だけで，medium，turbo，FTの上乗せの約半分〜9割が得られた。
ただし，包絡だけと融合の差が有意でないことは，両者が同等であることを示すものではない。
有声らしいフレームと巡回シフトしたフレームで差がなかったことから，効果が特定の区間に集まる証拠は得られなかった。
融合による変化は気導マイク音声への一様な接近では説明できず，近づく特徴と遠ざかる特徴があった。

包絡の成分は，系統の異なる認識器には逆に働いた。TAPSでは，複数のWhisperへの上乗せも，XLSRとZipformerの悪化も，包絡だけを加えたときに現れた。
一方VibraVoxでは，CE-SEと融合は仏語のXLSRとMMSも改善した。系統の異なる認識器への影響の正負はデータセットによって変わり，原因は特定できていない。
今回の2つのデータセットの間では，一方で適応した生成モデルは，適応し直さずには他方で使えなかった。ただし言語と機器が同時に変わっており，両者の寄与は分けていない。

限界として，今回は喉マイクごとに対データで適応しており，未知の機器への汎化は示せていない。VibraVoxでの評価も，話者は学習と重ならないが同一コーパス内の評価である。
評価は両データセットの収録音声のみで，喉マイクを使う主な理由である騒音下での効果は調べていない。
生成部分の処理時間は約1/10まで減らせたが，ステップ数を減らすとXLSRとZipformerのCERは悪化した。
また，元の学習と学習し直した学習とでは計算環境も異なり，系統の異なる認識器への悪化の大きさは学習ごとに変わった。\ref{sec:res_decomp}節の分解は元の学習についての結果である。
VibraVoxでの少量の対データの結果は1回の試行による。

%────────────────────────────────────────
\section{まとめ}
喉マイク音声の前処理が認識を助ける条件と，上乗せに関わる成分を，TAPSで9種，VibraVoxで7種の認識器を用いて調べた。
事前学習済みの生成モデルを対象の喉マイクで適応すると，FT Whisperを含む9種すべての認識器でCERが下がった。
今回の平均的な振幅特性と帯域制限による疑似喉マイクは，喉マイクの認識の難しさを再現しなかった。探索的な分析では，Whisper-smallで生成後も激音の字母の誤りが相対的に多く残った。
Whisperの損失で学習したSEとの融合は，Whisper-small，medium，turbo，FT Whisperに，3回の学習すべてで有意な小さな上乗せを足した。baseは3回中1回のみ有意で，Qwen3はどの回も有意でなかった。
元の学習では，融合と生成×4の差の包絡だけで，medium，turbo，FTの上乗せの約半分〜9割が得られた。
今回の2つのデータセットの間では，一方で適応した生成モデルは適応し直さずには他方で使えなかった。少量の対データでの適応と，未知の機器への汎化が今後の課題である。

\footnotesize
\let\oldbibitem\bibitem
\renewcommand{\bibitem}{\setlength{\itemsep}{0pt}\setlength{\parsep}{0pt}\oldbibitem}
\bibliographystyle{IEEEtran}
\bibliography{refs}

\end{document}
